
\setcounter{footnote}{42}

\SGARefTarget{sga3:iv:subsection:4:6}\subsection*{4.6. Description of the Quotient of a Sheaf by an Equivalence Relation}
\addcontentsline{toc}{subsection}{\SGARefLink{sga3:iv:subsection:4:6}{4.6}. Quotient of a Sheaf by an Equivalence Relation}

Recall that the given topology \(\mathcal T\) is assumed to be coarser than
the canonical topology.

\medskip
\noindent\SGARefTarget{sga3:iv:proposition:4:6:1}\textbf{Proposition 4.6.1.}
\textit{Let
\(
 R \mathrel{\substack{\xrightarrow{\ p_1\ }\\[-0.6ex]
                       \xrightarrow[\ p_2\ ]{}}} X
\)
be a \(\widetilde{\mathscr C}\)-equivalence relation in the sheaf \(X\).
Define \(F\in\operatorname{Ob}(\widehat{\mathscr C})\) as follows: for
each \(S\in\operatorname{Ob}(\mathscr C)\),\footnote{Editor's note:
\(R\times S\) denotes the equivalence relation in \(X\times S\) defined
by \(R\times S_{\mathrm{diag}}\subset
X\times X\times S\times S\), and \(R_Z\) denotes the equivalence
relation that it induces on \(Z\) (cf. \SGARefLink{sga3:iv:definition:3:1:6}{3.1.6}).}}
\[
 F(S)=\left\{\begin{array}{l}
 \text{sub-\(S\)-sheaves \(Z\) of \(X_S\), stable under
 \(R\times S\), whose quotient}\cr
 \text{by \(R_Z\) is \(S\), that is, for which
 \(R_Z\rightrightarrows Z\to S\) is exact}
\end{array}\right\}.
\]
\textit{For every sheaf \(Y\), \(\operatorname{Hom}(Y,F)\) is naturally
identified with the set}
\[
 \left\{\begin{array}{l}
 \text{sub-\(Y\)-sheaves of \(X\times Y\), stable under \(R\times Y\),}\cr
 \text{and having quotient \(Y\)}
 \end{array}\right\}.
\]
\textit{In particular, the subsheaf \(R\subset X\times X\) corresponds
to an element \(p\in\operatorname{Hom}(X,F)\), and the exact diagram}
\[
 R\mathrel{\substack{\xrightarrow{\ p_1\ }\\[-0.6ex]
                     \xrightarrow[\ p_2\ ]{}}}X\xrightarrow{p}F
\]
\textit{identifies \(F\) with the sheaf quotient \(X/R\).}

Indeed, put \(Q=X/R\). For every sheaf \(Y\) and every morphism
\(f\in\operatorname{Hom}(Y,Q)\), corresponding to a section
\(s:Y\hookrightarrow Q\times Y\), consider the diagram
\[
\tag{*}\label{sga3:iv:proposition:4:6:1:eq:unnamed:01}
\begin{array}{ccccc}
 R\times Y&\rightrightarrows&X\times Y&\longrightarrow&Q\times Y\cr
 &&\big\uparrow&&\big\uparrow\!{\scriptstyle s}\cr
 &&Z&\longrightarrow&Y,
\end{array}
\]
in which the right-hand square is cartesian. By \SGARefLink{sga3:iv:proposition:4:4:11}{4.4.11}, \(Z\) is
immediately seen to be a sub-\(Y\)-sheaf of \(X\times Y\), stable under
\(R\times Y\), with quotient \(Y\). Conversely, every \(Z\) of this
kind arises from a unique section of \(Q\times Y\) over \(Y\). Taking
\(Y\) first to be representable gives an isomorphism \(Q\simeq F\);
taking arbitrary \(Y\) then gives the asserted description of
\(\operatorname{Hom}(Y,F)\). Finally, the canonical morphism \(X\to Q\)
corresponds to the sub-\(X\)-sheaf \(R\subset X\times X\), which completes
the proof.

\medskip
\noindent\SGARefTarget{sga3:iv:corollary:4:6:2}\textbf{Corollary 4.6.2.}
\textit{Let \(G\) be any subfunctor of \(F\) such that
\(\operatorname{Hom}(X,G)\subset\operatorname{Hom}(X,F)\) contains
\(R\). Then the canonical morphism \(p:X\to F\) factors through \(G\).}

Since \(p\) is covering (\SGARefLink{sga3:iv:proposition:4:4:9}{4.4.9} and \SGARefLink{sga3:iv:proposition:4:4:3}{4.4.3}), it follows that \(G\) is a
refinement of \(F\). In particular, every subsheaf \(G\) of \(F\)
satisfying \SGARefLink{sga3:iv:corollary:4:6:2}{the preceding condition} equals \(F\) (\SGARefLink{sga3:iv:corollary:4:3:12}{4.3.12}).

\medskip
\noindent\SGARefTarget{sga3:iv:para:4:6:3}\textbf{4.6.3.}
We now consider the case in which \(X\) and \(R\) are representable.
We begin with some terminology. Besides the conditions \SGARefLink{sga3:iv:para:3:4:1:item:a:01}{(a)}--\SGARefLink{sga3:iv:para:3:4:1:item:d:01}{(d)}
introduced in \SGARefLink{sga3:iv:para:3:4:1}{3.4.1}, we shall use the following conditions on a family
\((M)\) of morphisms of \(\mathscr C\); for completeness, we repeat
conditions \SGARefLink{sga3:iv:para:4:6:3:item:a:01}{(a)}--\SGARefLink{sga3:iv:para:4:6:3:item:c:01}{(c)}:
\begin{enumerate}
 \item[(a)]\SGARefTarget{sga3:iv:para:4:6:3:item:a:01} \((M)\) is stable under base change.
 \item[(b)]\SGARefTarget{sga3:iv:para:4:6:3:item:b:01} The composite of two elements of \((M)\) belongs to \((M)\).
 \item[(c)]\SGARefTarget{sga3:iv:para:4:6:3:item:c:01} Every isomorphism belongs to \((M)\).
 \item[\((d_{\mathcal T})\)]\SGARefTarget{sga3:iv:para:4:6:3:item:d:mathcal:t:01} Every element of \((M)\) is covering.
 \footnote{Editor's note: \(\mathcal T\) denotes the given topology on
 \(\mathscr C\), which is coarser than the canonical topology.}
 \item[\((e_{\mathcal T})\)]\SGARefTarget{sga3:iv:para:4:6:3:item:e:mathcal:t:01} Let \(f:X\to Y\) be a morphism of
 \(\mathscr C\). If there is a refinement \(R\) of \(Y\) such that,
 for every \(Y'\to R\), the morphism
 \(X\times_Y Y'\to Y'\) belongs to \((M)\), then \(f\) belongs to
 \((M)\).
\end{enumerate}
Conditions \SGARefLink{sga3:iv:para:4:6:3:item:a:01}{(a)} and \SGARefLink{sga3:iv:para:4:6:3:item:b:01}{(b)} imply
\begin{enumerate}
 \item[\((a')\)]\SGARefTarget{sga3:iv:para:4:6:3:item:aprime:01} the fiber product of two elements of \((M)\) belongs
 to \((M)\).
\end{enumerate}
On the other hand, by \SGARefLink{sga3:iv:corollary:4:3:9}{4.3.9}, conditions \SGARefLink{sga3:iv:para:4:6:3:item:a:01}{(a)} and
\SGARefLink{sga3:iv:para:4:6:3:item:d:mathcal:t:01}{\((d_{\mathcal T})\)} imply
\begin{enumerate}
 \item[\((d')\)]\SGARefTarget{sga3:iv:para:4:6:3:item:dprime:01} every element of \((M)\) is a universal effective
 epimorphism.
\end{enumerate}

\medskip
\noindent\SGARefTarget{sga3:iv:para:4:6:4}\textbf{4.6.4.}
When \(\mathscr C\) has fiber products, the family of covering morphisms,
denoted \((M_{\mathcal T})\), satisfies the preceding conditions. Indeed
(cf. \SGARefLink{sga3:iv:definition:4:2:3}{4.2.3}), \SGARefLink{sga3:iv:para:4:6:3:item:a:01}{(a)} follows from \SGARefLink{sga3:iv:definition:4:2:3:item:c:1:01}{(C~1)}, \SGARefLink{sga3:iv:para:4:6:3:item:b:01}{(b)} from \SGARefLink{sga3:iv:definition:4:2:3:item:c:2:01}{(C~2)}, \SGARefLink{sga3:iv:para:4:6:3:item:c:01}{(c)} from \SGARefLink{sga3:iv:definition:4:2:3:item:c:4:01}{(C~4)},
\SGARefLink{sga3:iv:para:4:6:3:item:d:mathcal:t:01}{\((d_{\mathcal T})\)} from \SGARefLink{sga3:iv:definition:4:2:3}{the definition}, and
\SGARefLink{sga3:iv:para:4:6:3:item:e:mathcal:t:01}{\((e_{\mathcal T})\)} from \SGARefLink{sga3:iv:definition:4:2:3:item:c:5:01}{(C~5)}. Thus the results established below for
a family satisfying these conditions apply in particular to
\((M_{\mathcal T})\). In particular, \(\mathcal T\) may be taken to be
the canonical topology and \((M)\) to be the family of universal
effective epimorphisms.

\medskip
\noindent\SGARefTarget{sga3:iv:lemma:4:6:5}\textbf{Lemma 4.6.5.}
\textit{Let \((M)\) be a family of morphisms satisfying conditions
\SGARefLink{sga3:iv:para:4:6:3:item:a:01}{(a)}--\SGARefLink{sga3:iv:para:4:6:3:item:e:mathcal:t:01}{\((e_{\mathcal T})\)} above. Let \(R\) be a
\(\mathscr C\)-equivalence relation of type \((M)\) in
\(X\in\operatorname{Ob}(\mathscr C)\). Let \(\widetilde X\) be the
sheaf defined by \(X\), let \(\widetilde R\) be the
\(\widetilde{\mathscr C}\)-equivalence relation in \(\widetilde X\)
defined by \(R\), and let \(\widetilde X/\widetilde R\) be the sheaf
quotient. Then \(R\) is \((M)\)-effective if and only if
\(\widetilde X/\widetilde R\) is representable. When these conditions
hold, \(\widetilde X/\widetilde R\) is represented by the quotient
\(X/R\).}

Suppose first that \(R\) is \((M)\)-effective, and put \(Y=X/R\). The
canonical morphism \(p:X\to Y\) belongs to \((M)\), hence is covering by
\SGARefLink{sga3:iv:para:4:6:3:item:d:mathcal:t:01}{\((d_{\mathcal T})\)}. The corresponding morphism
\[
 \widetilde p:\widetilde X\longrightarrow
 \widetilde X/\widetilde R
\]
is therefore a universal effective epimorphism in
\(\widetilde{\mathscr C}\) (\SGARefLink{sga3:iv:proposition:4:4:3}{4.4.3}). Consequently it identifies
\(\widetilde X/\widetilde R\) with the quotient of \(\widetilde X\) by
the equivalence relation \(R'\) defined by \(\widetilde p\). Since the
canonical functor \(\mathscr C\to\widetilde{\mathscr C}\) commutes with
fiber products, \(R'\) is precisely \(\widetilde R\), because \(R\) is
the equivalence relation defined by \(p\).

Conversely, suppose that \(\widetilde X/\widetilde R\) is represented by
an object \(Y\) of \(\mathscr C\). Let \(p:X\to Y\) be the morphism
deduced from the canonical morphism
\(\widetilde X\to\widetilde X/\widetilde R\); by \SGARefLink{sga3:iv:proposition:4:4:3}{4.4.3} it is covering.
As before, \(R\) is plainly the equivalence relation defined by \(p\).
It remains only to show that \(p\in(M)\). The cartesian diagram
\[
\begin{array}{ccccc}
 R&\xrightarrow{\ \sim\ }&X\times_YX&\longrightarrow&X\cr
 {\scriptstyle p_2}\searrow&&\big\downarrow\!{\scriptstyle\mathrm{pr}_2}
 &&\big\downarrow\!{\scriptstyle p}\cr
 &&X&\xrightarrow{p}&Y
\end{array}
\]
shows that, after the covering base change \(p\), the morphism \(p\)
becomes \(p_2\), which belongs to \((M)\). \SGARefLink{sga3:iv:lemma:4:6:5}{The assertion} follows from
\SGARefLink{sga3:iv:para:4:6:3:item:e:mathcal:t:01}{\((e_{\mathcal T})\)}.

\medskip
\noindent\SGARefTarget{sga3:iv:corollary:4:6:5:1}\textbf{Corollary 4.6.5.1.}
\footnote{Editor's note: \SGARefLink{sga3:iv:corollary:4:6:5:1}{This corollary} has been added.}
\textit{Let \((M)\) satisfy conditions \SGARefLink{sga3:iv:para:4:6:3:item:a:01}{(a)}--\SGARefLink{sga3:iv:para:4:6:3:item:e:mathcal:t:01}{\((e_{\mathcal T})\)} above,
and let \(f:G\to G'\) be a morphism of \(\mathscr C\)-groups such that
\(f\in(M)\). Suppose that \(\operatorname{Ker}(f)\) is representable
(as is the case when \(\mathscr C\) has a final object \(e\)). Then the
equivalence relation in \(G\) defined by
\(H=\operatorname{Ker}(f)\) is \((M)\)-effective, and \(G'\) represents
the sheaf quotient \(\widetilde G/\widetilde H\) for the topology
\(\mathcal T\).}

This follows from \SGARefLink{sga3:iv:lemma:4:6:5}{4.6.5} and \SGARefLink{sga3:iv:corollary:3:4:7:1}{3.4.7.1}.

We can now state the main result of this subsection.

\medskip
\noindent\SGARefTarget{sga3:iv:theorem:4:6:6}\textbf{Theorem 4.6.6.}
\textit{Let \((M)\) be a family of morphisms satisfying Axioms
\SGARefLink{sga3:iv:para:4:6:3:item:a:01}{(a)}--\SGARefLink{sga3:iv:para:4:6:3:item:e:mathcal:t:01}{\((e_{\mathcal T})\)} of \SGARefLink{sga3:iv:para:4:6:3}{4.6.3}. Let \(R\) be a
\(\mathscr C\)-equivalence relation of type \((M)\) (cf. \SGARefLink{sga3:iv:definition:3:4:3}{3.4.3}) in an
object \(X\) of \(\mathscr C\). Consider the functor
\(F\in\operatorname{Ob}(\widehat{\mathscr C})\) defined by}
\[
 F(S)=\left\{\begin{array}{l}
 \text{sub-\(S\)-sheaves \(Z\) of \(X_S\), stable under \(R\times S\),}\cr
 \text{whose quotient by \(R_Z\) is \(S\)}
 \end{array}\right\}.
\]
\textit{Let \(F_0\) be the subfunctor of \(F\) defined as follows:
\(F_0(S)\) consists of the representable \(Z\in F(S)\), that is,}
\[
 F_0(S)=\left\{\begin{array}{l}
 \text{sub-\(\mathscr C/S\)-objects \(Z\) of \(X_S\), stable under
 \(R\times S\), such that}\cr
 \text{\(R_Z\) is \((M)\)-effective with quotient \(S\), i.e. such that}\cr
 \text{\(Z\to S\) belongs to \((M)\) and
 \(R_Z\simeq Z\times_S Z\)}
 \end{array}\right\}.
\]
\textit{Then:}
\begin{enumerate}
 \item[(i)]\SGARefTarget{sga3:iv:theorem:4:6:6:item:i:01} \textit{The morphism
 \(p\in\operatorname{Hom}(X,F)=F(X)\) defined by the subobject
 \(R\subset X\times X\) identifies \(F\) with the sheaf quotient of
 \(X\) by \(R\).}
 \item[(ii)]\SGARefTarget{sga3:iv:theorem:4:6:6:item:ii:01} \textit{The following conditions are equivalent:}
 \begin{enumerate}
  \item[(a)]\SGARefTarget{sga3:iv:theorem:4:6:6:item:a:01} \textit{\(F\) is representable;}
  \item[(b)]\SGARefTarget{sga3:iv:theorem:4:6:6:item:b:01} \textit{\(F_0\) is representable;}
  \item[(c)]\SGARefTarget{sga3:iv:theorem:4:6:6:item:c:01} \textit{\(R\) is \((M)\)-effective.}
 \end{enumerate}
 \textit{Under these conditions, \(F=F_0=X/R\).}
 \item[(iii)]\SGARefTarget{sga3:iv:theorem:4:6:6:item:iii:01} \textit{Let \((N)\) be a family of morphisms stable under
 base change, such that for every covering family \(\{S_i\to S\}\) and
 every family \(\{T_i\to S_i\}\) of morphisms in \((N)\), every
 descent datum on the \(T_i\) relative to \(\{S_i\to S\}\) is
 effective. Suppose that \(X\) is squarable (cf. \SGARefLink{sga3:iv:definition:1:6:0}{1.6.0}) and that
 \(R\to X\times X\) belongs to \((N)\). Then \(F_0=F\).}
\end{enumerate}

\noindent\textit{Proof.}
Assertion \SGARefLink{sga3:iv:theorem:4:6:6:item:i:01}{(i)} was proved in \SGARefLink{sga3:iv:proposition:4:6:1}{4.6.1}.

For \SGARefLink{sga3:iv:theorem:4:6:6:item:ii:01}{(ii)}, we have already proved the equivalence of \SGARefLink{sga3:iv:theorem:4:6:6:item:a:01}{(a)} and \SGARefLink{sga3:iv:theorem:4:6:6:item:c:01}{(c)}, as well
as the equality \(F=X/R\). It remains to prove that either \SGARefLink{sga3:iv:theorem:4:6:6:item:b:01}{(b)} or \SGARefLink{sga3:iv:theorem:4:6:6:item:c:01}{(c)}
implies \(F_0=F\). First note, as stated in \SGARefLink{sga3:iv:theorem:4:6:6}{the theorem}, that \(F_0\) is
indeed a subfunctor of \(F\): if \(S\in\operatorname{Ob}(\mathscr C)\)
and \(Z\in F_0(S)\), then \(Z\to S\) is squarable, and hence
\(Z\times_SS'\in F_0(S')\) for every \(S'\to S\). Since
\(R\in F(X)\) belongs to \(F_0(X)\), Corollary \SGARefLink{sga3:iv:corollary:4:6:2}{4.6.2} shows that \SGARefLink{sga3:iv:theorem:4:6:6:item:b:01}{(b)}
implies \(F_0=F\).

Now suppose that \SGARefLink{sga3:iv:theorem:4:6:6:item:c:01}{(c)} holds, and let \(Q\) be an object of \(\mathscr C\)
representing \(X/R\). Then \(X\to Q\) belongs to \((M)\). For every
\(S\in\operatorname{Ob}(\mathscr C)\) and every \(Z\in F(S)\), diagram
\(\SGARefLink{sga3:iv:proposition:4:6:1:eq:unnamed:01}{(*)}\) of \SGARefLink{sga3:iv:proposition:4:6:1}{4.6.1} gives
\[
 Z=S\times_{Q\times S}(X\times S).
\]
Thus \(Z\) is representable and \(Z\to S\) belongs to \((M)\), so
\(Z\in F_0(S)\).

For \SGARefLink{sga3:iv:theorem:4:6:6:item:iii:01}{(iii)}, let \(f\in\operatorname{Hom}(S,F)\) correspond to
\(Z\in F(S)\). We must show that \(f\) factors through \(F_0\), or
equivalently that \(Z\) is representable. This is immediate when \(f\)
factors through \(X\), by \SGARefLink{sga3:iv:lemma:4:6:7}{the following lemma}.

\medskip
\noindent\SGARefTarget{sga3:iv:lemma:4:6:7}\textbf{Lemma 4.6.7.}
\textit{Let \(x_0\in X(S)\). The image of \(x_0\) in \(F(S)\)
corresponds to the subsheaf \(Z\subset X_S\) defined by the two
cartesian squares}
\[
\begin{array}{ccccc}
 X_S&\xrightarrow{\mathrm{id}_{X_S}\times x_0}&
 X_S\times_SX_S&\longrightarrow&X\times X\cr
 \big\uparrow&&\big\uparrow&&\big\uparrow\cr
 Z&\longrightarrow&R_S&\longrightarrow&R.
\end{array}
\]

This follows at once from the description of the morphism \(X\to F\).

Returning to the proof of \SGARefLink{sga3:iv:theorem:4:6:6}{the theorem}, suppose first that \(f\) factors
through \(X\). Then \(Z\) is representable; moreover, since
\(R\to X\times X\) belongs to \((N)\), so does \(Z\to X_S\).

In general, \(f\) need not factor through \(X\). Since \(X\to F\) is
covering (\SGARefLink{sga3:iv:proposition:4:4:3}{4.4.3}), however, \SGARefLink{sga3:iv:remark:4:4:8:item:vii:01}{4.4.8(vii)} gives a covering family
\(\{S_i\to S\}\) and, for each \(i\), a morphism \(S_i\to X\) making
the diagram
\SGARefTarget{sga3:iv:supp:lemma:4:6:7:diagram:01}\[
\begin{array}{ccc}
 X&\longrightarrow&F\cr
 \big\uparrow&&\big\uparrow\!{\scriptstyle f}\cr
 S_i&\longrightarrow&S
\end{array}
\]
commute. By the preceding argument, the morphism
\(f_i:S_i\to F\) defined by \SGARefLink{sga3:iv:supp:lemma:4:6:7:diagram:01}{this diagram} belongs to
\(\operatorname{Hom}(S_i,F_0)\) and corresponds to the subsheaf
\(Z\times_SS_i\subset X_{S_i}\). The morphism
\(Z\times_SS_i\to X_{S_i}\) belongs to \((N)\), and the family
\(\{X_{S_i}\to X_S\}\) is covering. It therefore remains only to prove
\SGARefLink{sga3:iv:proposition:4:6:8}{the following proposition}.

\medskip
\noindent\SGARefTarget{sga3:iv:proposition:4:6:8}\textbf{Proposition 4.6.8.}
\textit{Let \(\{S_i\to S\}\) be a covering family and let \(Z\) be a
sheaf over \(S\). Suppose that, for every \(i\), the
\(S_i\)-functor \(Z\times_SS_i\) is represented by an object \(T_i\).
Then the family \((T_i)\) carries a canonical descent datum relative to
\(\{S_i\to S\}\). The sheaf \(Z\) is representable if and only if this
datum is effective; in that case the descended object represents \(Z\).}

First observe that, by \SGARefLink{sga3:iv:proposition:4:4:3}{4.4.3}, \(\{S_i\to S\}\) is a universal effective
epimorphic family in \(\widetilde{\mathscr C}\), and hence a descent
family in \(\widetilde{\mathscr C}\) (\SGARefLink{sga3:iv:proposition:2:3}{2.3}). If \(Z\) is represented by
an object \(T\), then \(T\times_SS_i\), regarded as a sheaf, is
isomorphic to \(Z\times_SS_i\). Thus the descent datum on the \(T_i\) is
effective, and its uniquely determined descended object is isomorphic to
\(T\), hence represents \(Z\).

Conversely, suppose that the canonical descent datum on the \(T_i\) is
effective, and let \(T\) be the descended object. Because
\(\{S_i\to S\}\) is a descent family in
\(\widetilde{\mathscr C}\), there is an \(S\)-morphism \(T\to Z\) whose
base change to each \(S_i\) is the canonical morphism
\(T_i\to Z\times_SS_i\). This morphism is locally an isomorphism; since
both \(T\) and \(Z\) are sheaves, \SGARefLink{sga3:iv:corollary:4:5:8}{4.5.8} implies that it is an
isomorphism.

\medskip
\noindent\SGARefTarget{sga3:iv:corollary:4:6:9}\textbf{Corollary 4.6.9.}
\textit{Let \(R\) be an \((M)\)-effective equivalence relation in
\(X\). For every sheaf \(F\), the map}
\[
 \operatorname{Hom}(X/R,F)\longrightarrow\operatorname{Hom}(X,F)
\]
\textit{identifies the first set with the subset of the second consisting
of morphisms compatible with \(R\).}

\medskip
\noindent\SGARefTarget{sga3:iv:corollary:4:6:10}\textbf{Corollary 4.6.10.}
\textit{Let \(\mathcal T'\) be a topology coarser than \(\mathcal T\),
for which the morphisms in \((M)\) are covering. Under the hypotheses of
\SGARefLink{sga3:iv:theorem:4:6:6:item:iii:01}{4.6.6(iii)}, \(X/R\) is also the sheaf quotient of \(X\) by \(R\) for
every topology lying between \(\mathcal T'\) and the canonical topology.}

\medskip
\noindent\SGARefTarget{sga3:iv:remark:4:6:11}\textbf{Remark 4.6.11.}
If, in the statement of \SGARefLink{sga3:iv:theorem:4:6:6:item:iii:01}{4.6.6(iii)}, one assumes in addition that the
descended morphism \(T\to S\) belongs to \((N)\) whenever the hypotheses
there are satisfied, then the inclusion morphisms \(Z\hookrightarrow X_S\)
also belong to \((N)\). This follows immediately from the construction of
\(Z\) by descent.

\medskip
\noindent\SGARefTarget{sga3:iv:remark:4:6:12}\textbf{Remark 4.6.12.}
The implications
\[
 \text{\SGARefLink{sga3:iv:theorem:4:6:6:item:c:01}{(c)}}\Longrightarrow\text{\SGARefLink{sga3:iv:theorem:4:6:6:item:b:01}{(b)}}\Longrightarrow\text{\SGARefLink{sga3:iv:theorem:4:6:6:item:a:01}{(a)}}
 \qquad\text{and}\qquad
 \text{\SGARefLink{sga3:iv:theorem:4:6:6:item:c:01}{(c)}}\Longrightarrow[F_0=F=X/R]
\]
were proved without using the sufficiency direction of Lemma \SGARefLink{sga3:iv:lemma:4:6:5}{4.6.5}, the
only place where condition \SGARefLink{sga3:iv:para:4:6:3:item:e:mathcal:t:01}{\((e_{\mathcal T})\)} enters. They therefore
remain valid when \((M)\) satisfies only conditions
\SGARefLink{sga3:iv:para:4:6:3:item:a:01}{(a)}--\SGARefLink{sga3:iv:para:4:6:3:item:d:mathcal:t:01}{\((d_{\mathcal T})\)}. The family of squarable covering morphisms is
an example of such a family \((M)\) (compare \SGARefLink{sga3:iv:para:4:6:4}{4.6.4}). For the canonical
topology these are precisely the universal effective epimorphisms. Hence:

\medskip
\noindent\SGARefTarget{sga3:iv:corollary:4:6:13}\textbf{Corollary 4.6.13.}
\textit{Let \(R\) be a universal effective equivalence relation in
\(X\). Then the object \(X/R\) of \(\mathscr C\) is the sheaf quotient
of \(X\) by \(R\) for the canonical topology. It represents the
following functor: \((X/R)(S)\) is the set of
sub-\(\mathscr C/S\)-objects \(Z\) of \(X_S\), stable under
\(R\times S\), for which the induced equivalence relation is universal
effective and has quotient \(S\).}

Similarly, for an arbitrary topology:

\medskip
\noindent\SGARefTarget{sga3:iv:corollary:4:6:14}\textbf{Corollary 4.6.14.}
\textit{Let \((M)\) be the family of squarable covering morphisms. If
\(R\) is an \((M)\)-effective equivalence relation in \(X\), then the
object \(X/R\) of \(\mathscr C\) is the sheaf quotient of \(X\) by
\(R\), and it represents the functor \(F_0\) of \SGARefLink{sga3:iv:theorem:4:6:6}{4.6.6}.}

\medskip
\noindent\SGARefTarget{sga3:iv:scholium:4:6:15}\textbf{Scholium 4.6.15.}
We can now make the following clarifications to \SGARefLink{sga3:iv:scholium:4:5:5}{4.5.5}. In questions
involving inverse limits alone---fiber products, algebraic structures,
and so forth---the results of Expos\'e I and \SGARefLink{sga3:iv:scholium:4:5:5}{4.5.5} allow us to identify
\(\mathscr C\), without distinction, with a full subcategory of either
\(\widetilde{\mathscr C}\) or \(\widehat{\mathscr C}\). The same is not
true for questions that mix inverse and direct limits.

Whenever both inverse and direct limits occur, and especially when
passing to quotients (for example, when putting a group structure on the
quotient of a group by an invariant subgroup), we shall regard the given
category as embedded in the category of sheaves. Thus, if \(R\) is a
\(\mathscr C\)-equivalence relation in an object \(X\) of
\(\mathscr C\), the notation \(X/R\) will mean the sheaf quotient of
\(X\) by \(R\), formerly denoted \(j(X)/j(R)\), and, when this sheaf is
representable, also the object that represents it. The preceding results
show that, in the most important cases, a quotient in \(\mathscr C\) is
also a quotient in the category of sheaves. In any event, we shall not
use the notation \(X/R\) for a quotient in \(\mathscr C\) that does not
coincide with the quotient in \(\widetilde{\mathscr C}\)---for instance,
one that is not universal. This modifies the definitions of \SGARefLink{sga3:iv:section:3}{no.~3}.

To study a problem of this kind, one therefore works first in the
category of sheaves, where all the usual results hold (cf. \SGARefLink{sga3:iv:subsection:4:4}{no.~4.4}), and
then specializes the resulting statements to the original category,
using the results of the present subsection and, when available,
criteria for the effectivity of descent. Examples of this method will
appear in the following subsections.
